\begin{tabularx}{\textwidth}{L{29mm}X X}
\toprule
\textbf{Modell} & \textbf{Alapgondolat} & \textbf{Mikor kedvező?} \\
\midrule
Vízesésmodell & A fő tevékenységeket különálló, egymást követő fázisokként kezeli. Egy fázis eredménye dokumentum, amelyre a következő fázis épít. & Ha a követelmények előre jól ismertek, stabilak, és a projekt jól dokumentált, menedzselhető folyamatot igényel. \\
Evolúciós / iteratív fejlesztés & A specifikáció, a fejlesztés és a validáció tevékenységei összefonódnak; a rendszer több verzión keresztül alakul. & Ha a megrendelő igényei nem teljesen tiszták, gyors visszajelzés kell, vagy a rendszer felhasználói tapasztalat alapján formálódik. \\
Komponensalapú fejlesztés & A rendszer lehetőség szerint meglévő, újrafelhasználható komponensekből épül fel; a követelmények részben igazodhatnak az elérhető komponensekhez. & Ha sok megbízható komponens áll rendelkezésre, és fontos a költség, a kockázat vagy a fejlesztési idő csökkentése. \\
Inkrementális fejlesztés & A rendszer szolgáltatásait kisebb, priorizált inkremensekben szállítják le, miközben a rendszer fokozatosan bővül. & Ha a megrendelőnek korán használható részrendszerre van szüksége, és a funkciók jól priorizálhatók. \\
Spirális fejlesztés & A fejlesztést ismétlődő ciklusokként ábrázolja, amelyekben hangsúlyos a célok kijelölése, a kockázatelemzés, a fejlesztés-validálás és a következő ciklus tervezése. & Nagyobb, kockázatosabb rendszereknél, ahol a kockázatok tudatos kezelése különösen fontos. \\
\bottomrule
\end{tabularx}
